\chapter{Background}\label{ch:Background}

The models analyzed in this thesis draw on a shared set of 
technical foundations: the structure of SVG itself, the 
representations used to manipulate it during learning, the 
rendering bridge that connects vector parameters to pixel 
space, and the supervision signals that drive optimization. 
Understanding these components is what makes it possible to 
compare models not just by their outputs, but by the 
structural choices that produce them.


\section{Vector Graphics Fundamentals}

Scalable Vector Graphics (SVG) is an XML-based standard for 
representing two-dimensional graphics as structured 
text~\cite{eisenbergSVGEssentialsProducing20}. Unlike raster formats, an SVG 
file is a readable document that describes how an image should 
be drawn rather than storing pixel values. At its core, an SVG 
file consists of instructions that define geometric primitives, 
styling, and hierarchical relationships between 
elements~\cite{quintScalableVectorGraphics2003}.

SVG supports a range of graphic elements, including basic shapes 
such as lines, rectangles, and circles, as well as text and 
raster image embeds. The most expressive element is the path, 
which allows complex shapes to be defined using sequences of 
commands: moving a virtual pen, drawing straight lines, and 
creating smooth curves~\cite{eisenbergSVGEssentialsProducing20}.

SVG also supports grouping and layered organization of elements. 
Multiple shapes can be combined into groups that behave as a 
single unit for styling and editing~\cite{quintScalableVectorGraphics2003}. 
Elements are rendered in a defined order, and later elements 
appear on top of earlier ones, making layering an intrinsic 
part of the representation. This hierarchical structure matters 
for generation: models that produce only geometric approximations 
without respecting layer organization generate outputs that are 
visually plausible but difficult to edit.

\subsection{Paths and Bézier Curves}

Most SVG graphics are built using paths. A path is a sequence 
of commands --- \emph{move to}, \emph{line to}, \emph{curve to} 
--- that define its shape~\cite{eisenbergSVGEssentialsProducing20}. Curved 
segments use Bézier curves, which are smooth curves controlled 
by a set of control points that influence curvature without 
lying directly on the curve.

SVG supports both quadratic Bézier curves, which use a single 
control point, and cubic Bézier curves, which use two and offer 
greater flexibility for complex shapes~\cite{eisenbergSVGEssentialsProducing20}. 
From the perspective of SVG generation, these control points 
are the core geometric parameters that most models either 
predict directly or optimize iteratively.

\subsection{Coordinate Systems and Transformations}

SVG graphics are drawn within a coordinate system defined by a 
viewport, with the origin at the top-left corner, the x-axis 
pointing right, and the y-axis pointing 
downward~\cite{quintScalableVectorGraphics2003}. Transformations --- 
translation, scaling, rotation, skewing --- can be applied to 
individual shapes or to groups, allowing collections of paths 
to be moved or resized while preserving their internal 
structure. Nested groups define local coordinate frames, which 
is particularly relevant for layered compositions where 
subgroups must maintain their relative geometry independently.

\subsection{Scope of SVG in This Thesis}

SVG includes additional features such as clipping, masking, 
and filter effects. These primarily affect rendering appearance 
rather than geometric structure, and are not considered further 
here. This thesis treats SVG as a structured representation of 
geometric primitives, paths, and hierarchical relationships. 
Understanding these components is what allows the generation 
models analyzed in Chapters~\ref{ch:Landscape} 
through~\ref{ch:synthesis} to be compared on structural 
grounds --- not just by their visual output.



\section{Vector Representations}

SVG generation and vectorization models do not all manipulate 
vector graphics the same way internally. The final output is 
always an SVG file --- text-based, editable, conforming to the 
same standard. What differs is the internal representation 
used during learning or optimization, and this choice directly 
affects which algorithms can be applied, how stable training 
is, and what structural properties tend to emerge. The 
following subsections introduce the four representation types 
that appear across the models analyzed in this thesis.

\subsection{Parametric Representation}

SVG files are already parametric by nature: every element is 
defined by a set of numeric values --- control point 
coordinates, stroke widths, colors --- that can be read, 
modified, and written as plain text~\cite{eisenbergSVGEssentialsProducing20}. 
Parametric representations in generation models follow this 
same structure directly. A graphic is treated as a collection 
of primitives, each fully described by continuous parameters 
that can be adjusted independently. In optimization-based 
methods, generation works by initializing a set of shapes and 
iteratively refining their parameters until the rendered 
output matches a target~\cite{maLayerwiseImageVectorization2022}.

\subsection{Sequential Representation}

A different way to think about vector graphics is as a 
sequence of drawing instructions: move the pen here, draw a 
line there, curve to this point. This mirrors how SVG path 
data is actually written, but it also opens up a connection 
to language modelling --- a vector graphic encoded as a 
command sequence looks structurally similar to a 
sentence~\cite{carlierDeepSVGHierarchicalGenerative2020}. 
Command order determines both drawing sequence and layering, 
and autoregressive models can generate graphics token by 
token the same way they generate text. The trade-off is that 
errors accumulate: a mistake early in a long sequence can 
corrupt everything that follows.

\subsection{Implicit Neural Representation}

Implicit representations take a fundamentally different 
approach: there are no explicit shapes, no stored control 
points, no command sequences. Instead, a neural network 
defines the graphic entirely through its 
weights~\cite{thamizharasanNIVeLNeuralImplicit2024}. Given a 
spatial coordinate as input, the network outputs properties 
at that location --- color, layer assignment, shape 
membership. The geometry is never stored; it emerges from 
evaluating this continuous function across space. Because 
the entire representation is differentiable by construction, 
implicit models integrate naturally with gradient-based 
optimization and diffusion-based supervision.

\subsection{Latent Representation}

Latent representations are not vector representations in the 
geometric sense. Rather than describing the structure of a 
single graphic, a latent space compresses many graphics into 
a shared numerical embedding learned during training. A VAE, 
for example, encodes inputs into compact vectors and learns 
to reconstruct them through a 
decoder~\cite{kingmaAutoEncodingVariationalBayes2022}. The 
actual geometry --- paths, control points, commands --- still 
lives in one of the other representation types underneath. 
This concept is included here because several models in 
Chapter~\ref{ch:representation} use latent spaces as an 
organizational layer on top of their base representation, 
and that distinction matters when comparing how their 
outputs are structured.

\section{Differentiable Rendering}

Before neural networks can optimize vector graphics using 
image-based loss functions, there needs to be a way to connect 
the two spaces. Vector parameters --- control points, stroke 
widths, colors --- live in a continuous geometric space. Loss 
functions operate on pixel values. These two spaces are not 
naturally connected: traditional rasterization converts vector 
geometry into pixels, but this process is not differentiable, 
meaning gradients cannot flow back through it to update the 
vector parameters.

DiffVG, introduced by Li et al.~\cite{liDifferentiableVectorGraphics2020a}, 
solves this by approximating rasterization in a way that 
supports gradient computation. Given a set of vector 
parameters as input, DiffVG renders a pixel image as output. 
Because the rendering step is differentiable, a loss computed 
on the output image --- whether a pixel reconstruction loss, 
a perceptual loss, or a signal from a pretrained semantic 
model --- can be backpropagated all the way through to the 
vector parameters themselves. The system learns which control 
point to move, and in which direction, to reduce the error. 
Technically, DiffVG handles the discrete visibility changes 
and anti-aliasing effects that make standard rasterization 
non-differentiable~\cite{liDifferentiableVectorGraphics2020a}.

Before DiffVG, optimizing vector graphics with gradient-based 
methods was not straightforward. Its introduction made it 
possible to apply the full range of image-based supervision 
strategies --- reconstruction losses, CLIP guidance, diffusion 
priors --- directly to vector parameters. Most of the 
optimization-based models analyzed in this thesis rely on 
this capability, which is why it is introduced here as a 
foundational component rather than a model-specific detail.


\section{Loss Functions and Semantic Guidance}

Even when a model manipulates vector parameters internally, 
comparing outputs in pixel space is more practical than 
comparing vector structures directly. Most visual data exists 
as raster images, and pixel-based comparison gives access to 
well-established loss functions from computer vision. When 
combined with differentiable rendering, gradients from these 
losses flow back through the rasterizer to update vector 
parameters. This section introduces the main supervision 
signals used across the models analyzed in this thesis.

\subsection{Reconstruction Losses}

The simplest losses operate directly on pixel values. The L2 
loss measures the mean squared difference between predicted 
and target pixels, while the L1 loss measures the mean 
absolute difference~\cite{goodfellow2016deep}. Both are 
computationally efficient and provide stable gradients when 
a reference image is available. Their limitation is that 
they treat each pixel independently, without capturing 
higher-level structure such as shape boundaries or object 
regions.

Perceptual losses address this by comparing feature 
representations extracted from pretrained convolutional 
networks rather than raw pixel values. The formulation 
proposed by Johnson et al.~\cite{johnsonPerceptualLossesRealTime2016} 
uses intermediate feature maps from VGG~\cite{simonyan2015very}, 
which capture edges, textures, and object-level structure. 
Comparing at the feature level produces supervision signals 
that better reflect visual similarity as humans perceive it, 
and helps preserve overall structure even when exact pixel 
alignment differs slightly.

For tasks where shape boundary accuracy matters more than 
color fidelity, overlap-based losses are used. Intersection 
over Union (IoU)~\cite{everingham2010pascal} measures the 
overlap between predicted and target regions. In vector 
graphics, mask-based supervision of this kind encourages 
correct topology and spatial consistency across layers, 
rather than minimizing pixel intensity differences.

All three loss types depend on having a reference raster 
image. In text-to-SVG settings, where no target image 
exists, alternative guidance strategies are needed.

\subsection{CLIP}

CLIP (Contrastive Language--Image Pretraining), introduced 
by Radford et al.~\cite{radfordLearningTransferableVisual2021}, 
is trained on large-scale image--text pairs to learn a 
shared embedding space where semantically matching pairs 
are placed close together and mismatched pairs are pushed 
apart. Given an image and a text prompt, CLIP measures 
their similarity through separate image and text encoders, 
typically using cosine similarity between the resulting 
embeddings.

In SVG generation pipelines, the rendered vector graphic 
is passed through the CLIP image encoder while the text 
prompt is encoded separately. If the rendered output does 
not match the prompt semantically, the similarity score 
will be low. Because CLIP is differentiable with respect 
to its image input, and because images can be produced 
from vector parameters via differentiable rendering, the 
CLIP similarity score can serve as an optimization signal. 
Gradients flow from the similarity score back through the 
rasterizer to the SVG parameters, allowing the system to 
adjust shapes and colors toward better semantic alignment. 
CLIP does not generate images in this setup --- it acts as 
a semantic evaluator that provides a direction for 
improvement.

\subsection{Diffusion Models}

Diffusion models are generative models that learn to 
reverse a gradual noise corruption process. During 
training, clean images are progressively corrupted by 
adding noise, and the model learns to denoise them step 
by step. After training, new images can be generated by 
starting from random noise and iteratively 
denoising~\cite{ho2020denoising}. Stable 
Diffusion~\cite{rombachHighResolutionImageSynthesis2021} extends this by 
performing the diffusion process in a compressed latent 
space rather than directly in pixel space, significantly 
reducing computational cost while preserving image quality. 
When conditioned on text, diffusion models learn to 
generate images consistent with the textual description, 
capturing complex visual distributions and rich semantic 
structure.

\subsection{Score Distillation Sampling and VSD}

Score Distillation Sampling (SDS), introduced by 
DreamFusion~\cite{pooleDreamFusionTextto3DUsing2022} for 
text-to-3D generation and later adopted for vector 
graphics by VectorFusion~\cite{jainVectorFusionTexttoSVGAbstracting2022}, 
uses a pretrained text-conditioned diffusion model as an 
optimization signal. The current SVG is rendered to an 
image, noise is added according to the diffusion schedule, 
and the diffusion model predicts the denoised version 
conditioned on the text prompt. The difference between 
the noisy image and the predicted denoising defines a 
gradient, which is backpropagated through the rasterizer 
to update the SVG parameters. No ground-truth image is 
required --- the diffusion model's learned prior acts as 
the teacher.

A known limitation of SDS is oversaturation: because the 
gradient update treats the rendered image as if it were 
pure noise, results tend to be over-smoothed and lack 
fine detail. Variational Score Distillation (VSD), 
introduced in SVGDreamer~\cite{xingSVGDreamerTextGuided2024}, 
addresses this by modeling the rendered output as a 
particle in a distribution rather than treating it as 
noise directly. This produces more stable gradients and 
visually consistent results. Both SDS and VSD are used 
across several models analyzed in 
Chapter~\ref{ch:supervision}.



